\documentclass[UTF8]{ctexart}
\usepackage{graphicx} % Required for inserting images


设 \(f:\mathbb R^d\to\mathbb R\) 是 \(L\)-smooth，并且每次只能获得带噪函数值

\[
F(x)=f(x)+\xi,
\qquad
\mathbb E[\xi\mid x]=0,
\qquad
\mathbb E[\xi^2\mid x]\leq \gamma^2.
\]

考虑 \(m\) 个独立 Gaussian 方向的 antithetic ES estimator：

\[
\widehat g_\sigma(x)
=
\frac1m\sum_{i=1}^m
\frac{F(x+\sigma u_i)-F(x-\sigma u_i)}
{2\sigma}u_i,
\qquad
u_i\sim\mathcal N(0,I_d).
\]

假设正负两个函数值的噪声相互独立，则

\[
\boxed{
\begin{aligned}
\mathbb E\left\|
\widehat g_\sigma(x)-\nabla f(x)
\right\|^2
\leq\;&
L^2\sigma^2d\\
&+
\frac{2(d+1)}{m}
\|\nabla f(x)\|^2\\
&+
\frac{L^2\sigma^2}{2m}
d(d+2)(d+4)\\
&+
\frac{\gamma^2d}{2m\sigma^2}.
\end{aligned}
}
\]

其中最后一项就是函数值噪声造成的误差：

\[
\boxed{
\mathbb E\left\|
\widehat g_\sigma^{\rm noise}(x)
\right\|^2
\leq
\frac{\gamma^2d}{2m\sigma^2}.
}
\]

因此简化后可以写成

\[
\boxed{
\mathbb E\left\|
\widehat g_\sigma(x)-\nabla f(x)
\right\|^2
=
O\left(
L^2\sigma^2d
+
\frac{d}{m}\|\nabla f(x)\|^2
+
\frac{L^2\sigma^2d^3}{m}
+
\frac{\gamma^2d}{m\sigma^2}
\right).
}
\]

对应的均方根误差为

\[
\boxed{
\sqrt{
\mathbb E\left\|
\widehat g_\sigma(x)-\nabla f(x)
\right\|^2
}
=
O\left(
L\sigma\sqrt d
+
\|\nabla f(x)\|\sqrt{\frac dm}
+
\frac{L\sigma d^{3/2}}{\sqrt m}
+
\frac{\gamma}{\sigma}\sqrt{\frac dm}
\right).
}
\]

只平衡主要的平滑误差和函数值噪声：

\[
L^2\sigma^2d
\approx
\frac{\gamma^2d}{m\sigma^2},
\]

得到

\[
\boxed{
\sigma
\asymp
\sqrt{\frac{\gamma}{L\sqrt m}}
=
\left(\frac{\gamma^2}{mL^2}\right)^{1/4}.
}
\]

若考虑有限方向产生的 \(L^2\sigma^2d^3/m\) 项，则可以更完整地取

\[
\boxed{
\sigma^4
\asymp
\frac{\gamma^2d/(m)}
{L^2d+L^2d^3/m}
=
\frac{\gamma^2}
{L^2(m+d^2)},
}
\]

即

\[
\boxed{
\sigma
\asymp
\sqrt{\frac{\gamma}{L\sqrt{m+d^2}}}.
}
\]

这里最关键的是噪声项随 \(\sigma\to0\) 按照

\[
\frac{1}{\sigma^2}
\]

发散，因此带噪函数值下不能令 \(\sigma\) 任意趋近于零。\\

\newtheorem{proposition}{命题}
\newtheorem{corollary}{推论}

\title{带噪函数值下 ES 梯度估计的误差与 $\sigma$ 的选择}
\author{}
\date{}

\begin{document}
\maketitle

\section{背景}

考虑带有噪声函数评估的 antithetic ES 梯度估计器，对于 $f:\mathbb{R}^d\to\mathbb{R}$，我们只能观测到 $F(x) = f(x) + \xi$，其中 $\mathbb{E}[\xi\mid x]=0$ 且 $\mathbb{E}[\xi^2\mid x]\leq\gamma^2$：
\[
\widehat{g}_\sigma(x) = \frac{1}{m}\sum_{i=1}^m \frac{F(x+\sigma u_i) - F(x-\sigma u_i)}{2\sigma}u_i,
\quad u_i\sim\mathcal{N}(0,I_d).
\]

均方误差的上界为：
\[
\mathbb{E}\|\widehat{g}_\sigma(x) - \nabla f(x)\|^2
\leq
L^2\sigma^2 d + \frac{d}{m}\|\nabla f(x)\|^2 + \frac{L^2\sigma^2 d^3}{m} + \frac{\gamma^2 d}{m\sigma^2}.
\]

\section{主要结果}

\subsection{最小化绝对误差：$\sigma$ 与 $\|\nabla f\|$ 无关}

\begin{proposition}[绝对误差最小化]
\label{prop:abs}
最小化 $\mathbb{E}\|\widehat{g}_\sigma - \nabla f\|^2$ 的最优 $\sigma$ 为
\[
\sigma^* \propto \left(\frac{\gamma^2}{mL^2}\right)^{1/4},
\]
它与 $\|\nabla f(x)\|$ \textcolor{red}{无关}。
\end{proposition}

\begin{proof}
在四项误差中，只有第一项和最后一项依赖于 $\sigma$。平衡主导项：
\[
L^2\sigma^2 d \approx \frac{\gamma^2 d}{m\sigma^2}
\implies
L^2\sigma^4 \approx \frac{\gamma^2}{m}
\implies
\sigma^* \propto \left(\frac{\gamma^2}{mL^2}\right)^{1/4}.
\]
注意 $\|\nabla f(x)\|$ 只出现在第二项 $\frac{d}{m}\|\nabla f\|^2$ 中，而该项不依赖于 $\sigma$。
\end{proof}

\section{$\rho$ 作为可观测的信噪比}

前面的误差界依赖于 $\gamma$、$L$ 与 $\|\nabla f(x)\|$，而这三个量在实际的 LLM 微调中都无法直接测量。本节说明：训练过程中逐步记录的 split-half 秩相关系数 $\rho$ 并不是一个独立的诊断量，它恰好是单方向信噪比 $\sigma\|\nabla f\|/\gamma$ 的一个单调重参数化。因此可以把 $\rho$ 代入误差界的右端，消去 $\gamma$ 与 $\sigma$，得到一个只含 $d$、$m$、$\rho$ 的相对误差界。

\subsection{设置与记号}

第 $i$ 个 member 位于 $x + \sigma u_i$，其真实适应度记为 $f_i = f(x + \sigma u_i)$。实际观测到的是它在 $B$ 道题上的经验准确率
\[
\widehat f_i = f_i + \xi_i,
\qquad
\mathbb E[\xi_i] = 0,
\qquad
\operatorname{Var}(\xi_i) = \gamma^2 .
\]
注意 $\gamma^2 = \gamma_1^2 / B$，其中 $\gamma_1^2$ 是单题噪声方差；也就是说 $\gamma$ 依赖于批量大小 $B$。

定义跨 member 的两个方差：
\[
S^2 := \operatorname{Var}_i(f_i),
\qquad
V_{\mathrm{obs}} := \operatorname{Var}_i(\widehat f_i) = S^2 + \gamma^2 .
\]
$S^2$ 是 ES 更新真正依赖的\emph{信号}，$V_{\mathrm{obs}}$ 是可观测的 member 适应度方差。

split-half 相关系数 $\rho$ 的定义是：把 $B$ 道题随机劈成两个不相交的一半，每个 member 在两半上分别评分，再对两组得分求相关系数。

\subsection{桥梁：$\rho$ 与 $(\gamma,\sigma,\|\nabla f\|)$ 的关系}

\begin{proposition}[信号方差的一阶展开]
\label{prop:signal}
member 之间适应度的离散程度在一阶上就是 $\sigma\|\nabla f\|$：
\[
\boxed{
S^2 = \sigma^2\|\nabla f(x)\|^2 + O(\sigma^4 L^2 d),
\qquad
\frac{S}{\sigma} \leq \|\nabla f(x)\| + \sigma L\sqrt{d/2}.
}
\]
\end{proposition}

\begin{proof}
对 $u\sim\mathcal N(0,I_d)$ 做 Taylor 展开：
\[
f(x+\sigma u) - f(x)
=
\sigma\nabla f(x)^\top u
+
\frac{\sigma^2}{2}u^\top H u
+
\cdots,
\qquad H := \nabla^2 f(x).
\]
第一项的方差为 $\sigma^2\|\nabla f\|^2$。第二项的方差为
\[
\frac{\sigma^4}{4}\operatorname{Var}(u^\top H u)
=
\frac{\sigma^4}{4}\cdot 2\|H\|_F^2
=
\frac{\sigma^4\|H\|_F^2}{2}
\leq
\frac{\sigma^4 d L^2}{2},
\]
这里用了 $L$-smooth 蕴含 $\|H\|_2\leq L$，故 $\|H\|_F^2\leq d\|H\|_2^2\leq dL^2$。由标准差的三角不等式（Minkowski）
\[
S \leq \sigma\|\nabla f\| + \sigma^2 L\sqrt{d/2},
\]
两边除以 $\sigma$ 即得。
\end{proof}

\begin{proposition}[split-half 分解]
\label{prop:rho}
设两半各含 $B/2$ 道题且题目不相交，则
\[
\boxed{
\rho
=
\frac{S^2}{S^2 + 2\gamma^2}.
}
\]
\end{proposition}

\begin{proof}
每一半只用 $B/2$ 道题，故该半上的噪声方差为 $\gamma_1^2/(B/2) = 2\gamma^2$。两半的题目不相交，因此在给定 member 的条件下两半的噪声独立。记两半的得分为
\[
\widehat f_i^{(1)} = f_i + \eta_i^{(1)},
\qquad
\widehat f_i^{(2)} = f_i + \eta_i^{(2)},
\qquad
\operatorname{Var}(\eta_i^{(h)}) = 2\gamma^2,
\quad
\eta^{(1)}\perp\eta^{(2)}.
\]
于是跨 member 的协方差与方差分别为
\[
\operatorname{Cov}\big(\widehat f^{(1)},\widehat f^{(2)}\big) = S^2,
\qquad
\operatorname{Var}\big(\widehat f^{(h)}\big) = S^2 + 2\gamma^2,
\]
相关系数即为所述。
\end{proof}

\begin{corollary}[关键代换]
\label{cor:sub}
由命题 \ref{prop:rho} 反解并代入命题 \ref{prop:signal}：
\[
\boxed{
\frac{\gamma}{\sigma}
=
\frac{S}{\sigma}\sqrt{\frac{1-\rho}{2\rho}}
\;\leq\;
\left(\|\nabla f(x)\| + \sigma L\sqrt{d/2}\right)
\sqrt{\frac{1-\rho}{2\rho}},
}
\]
一阶下即 $\dfrac{\gamma}{\sigma}\approx\|\nabla f(x)\|\sqrt{\dfrac{1-\rho}{2\rho}}$。
\end{corollary}

\subsection{代入误差界}

\begin{proposition}[$\rho$-参数化的梯度估计误差]
\label{prop:bound}
把推论 \ref{cor:sub} 代入第 1 节的均方根误差界，$\gamma$ 与 $\sigma$ 从所有方差项中消失：
\[
\boxed{
\sqrt{\mathbb E\left\|\widehat g_\sigma(x)-\nabla f(x)\right\|^2}
=
O\left(
\underbrace{L\sigma\sqrt d + \frac{L\sigma d^{3/2}}{\sqrt m}}_{\text{平滑偏差，与 }\rho\text{ 无关}}
+
\underbrace{\|\nabla f(x)\|\sqrt{\frac dm}\left(1+\sqrt{\frac{1-\rho}{2\rho}}\right)}_{\text{全部方差项}}
\right).
}
\]
相应地，方差部分的\emph{相对}误差是一个只含 $d,m,\rho$ 的量：
\[
\boxed{
\frac{1}{\|\nabla f(x)\|}
\sqrt{\mathbb E\left\|\widehat g_\sigma^{\mathrm{var}}(x)\right\|^2}
=
\sqrt{\frac dm}\left(1+\sqrt{\frac{1-\rho}{2\rho}}\right).
}
\]
\end{proposition}

\begin{proof}
原界中只有第四项 $\frac{\gamma^2 d}{m\sigma^2}$ 含 $\gamma$。由推论 \ref{cor:sub}，
\[
\frac{\gamma}{\sigma}\sqrt{\frac dm}
\leq
\|\nabla f\|\sqrt{\frac dm}\sqrt{\frac{1-\rho}{2\rho}}
+
\sigma L\sqrt{\frac d2}\sqrt{\frac dm}\sqrt{\frac{1-\rho}{2\rho}} .
\]
第一部分与原界第二项 $\|\nabla f\|\sqrt{d/m}$ 同类，合并得到括号中的因子 $1+\sqrt{(1-\rho)/(2\rho)}$。第二部分为
\[
O\!\left(\frac{\sigma L d}{\sqrt m}\sqrt{\frac{1-\rho}{2\rho}}\right),
\]
当 $\frac{1-\rho}{2\rho}\lesssim d$（即 $\rho\gtrsim 1/d$）时它被原界中已有的 $\frac{L\sigma d^{3/2}}{\sqrt m}$ 吸收，可以略去。除以 $\|\nabla f\|$ 即得相对误差式。
\end{proof}

\subsection{三个推论}

\begin{corollary}[$\rho$ 即单方向信噪比]
\label{cor:snr}
沿单个方向 $u$ 的有限差分中，信号部分 $\frac{f(x+\sigma u)-f(x-\sigma u)}{2\sigma}\approx\nabla f(x)^\top u$ 的量级为 $\|\nabla f(x)\|$，噪声部分 $\frac{\xi_+-\xi_-}{2\sigma}$ 的标准差为 $\frac{\sqrt2\gamma}{2\sigma}\approx\frac{\gamma}{\sigma}$，故定义单方向信噪比
\[
\mathrm{SNR}_{\mathrm{single}} := \frac{\sigma\|\nabla f(x)\|}{\gamma}.
\]
由推论 \ref{cor:sub} 立得
\[
\boxed{
\mathrm{SNR}_{\mathrm{single}}
=
\sqrt{\frac{2\rho}{1-\rho}}.
}
\]
即 $\rho$ 与理论中的单方向信噪比是同一个量的单调变换；$\rho\to 0$ 当且仅当 $\mathrm{SNR}_{\mathrm{single}}\to 0$。这不是类比，而是恒等式：$\rho$ 是 $\mathrm{SNR}_{\mathrm{single}}$ 的\textcolor{red}{可观测替代量}。
\end{corollary}

\begin{corollary}[瓶颈判据：$\rho = 1/3$]
\label{cor:third}
命题 \ref{prop:bound} 的方差项由两部分构成：ES 固有的有限方向采样噪声（系数 $1$）与函数值评测噪声（系数 $\sqrt{(1-\rho)/(2\rho)}$）。两者相等当且仅当
\[
\frac{1-\rho}{2\rho} = 1
\quad\Longleftrightarrow\quad
\boxed{\rho = \tfrac13 .}
\]
因此：$\rho < 1/3$ 时评测噪声是瓶颈，应增大 $B$（或降低生成随机性、提高 reward 分辨率），增大 $m$ 收效甚微；$\rho > 1/3$ 时 population size 是瓶颈，应增大 $m$，增大 $B$ 是浪费算力。
\end{corollary}

\begin{corollary}[每步进展正比于 $\rho$]
\label{cor:progress}
在最优步长下，单步下降量为
\[
\Delta f^*
=
\frac{\|\nabla f\|^4}{2L\,\mathbb E\|\widehat g_\sigma\|^2}
=
\frac{\|\nabla f\|^2}{2L}
\cdot
\frac{1}{1 + \frac dm\cdot\frac{1+\rho}{2\rho}}
\;\xrightarrow[\ \rho\ll 1\ ]{}\;
\boxed{\frac{m\rho}{d}\cdot\frac{\|\nabla f\|^2}{L}.}
\]
即在小 $\rho$ 区，\textcolor{red}{$\rho$ 与 $m$ 完全可互换}：$\rho$ 减半等价于 population size 减半，而 $\rho = 0$ 意味着期望进展为零。
\end{corollary}

\begin{proof}
以 $x' = x-\alpha\widehat g_\sigma$ 作一步更新，由 $L$-smooth 与 $\mathbb E[\widehat g_\sigma]\approx\nabla f$，
\[
\mathbb E[f(x')] \leq f(x) - \alpha\|\nabla f\|^2 + \frac{L\alpha^2}{2}\mathbb E\|\widehat g_\sigma\|^2 ,
\]
对 $\alpha$ 最小化得 $\alpha^* = \|\nabla f\|^2/(L\,\mathbb E\|\widehat g_\sigma\|^2)$ 与 $\Delta f^* = \|\nabla f\|^4/(2L\,\mathbb E\|\widehat g_\sigma\|^2)$。再由推论 \ref{cor:sub}，
\[
\mathbb E\|\widehat g_\sigma\|^2
\approx
\|\nabla f\|^2 + \frac dm\|\nabla f\|^2 + \frac{\gamma^2 d}{m\sigma^2}
=
\|\nabla f\|^2\left[1 + \frac dm\left(1 + \frac{1-\rho}{2\rho}\right)\right]
=
\|\nabla f\|^2\left[1 + \frac dm\cdot\frac{1+\rho}{2\rho}\right].
\]
代入即得；$\rho\ll1$ 时分母 $\approx \frac{d}{2m\rho}$，给出末式。
\end{proof}

\subsection{全可观测形式}

由命题 \ref{prop:rho} 与 Spearman--Brown 公式，全长（$B$ 题）的 reliability 为
\[
R := \frac{2\rho}{1+\rho} = \frac{S^2}{S^2+\gamma^2},
\]
于是
\[
\boxed{
\widehat{\|\nabla f\|} = \frac{\sqrt{R\,V_{\mathrm{obs}}}}{\sigma},
\qquad
\widehat\gamma = \sqrt{(1-R)\,V_{\mathrm{obs}}} .
}
\]
$V_{\mathrm{obs}}$ 与 $\rho$ 都是训练过程中逐步可记录的量，因此 $\|\nabla f\|$ 与 $\gamma$ 的逐步曲线可以直接从已有的训练日志反推，无需额外计算。这使得"$\sigma$ 应当如何随 $\|\nabla f\|$ 变化"这一问题可以在实验上直接检验，而不必依赖对 $\gamma$ 与 $L$ 的先验假设。

\subsection{适用范围}

以上结论有四点限制，需要在使用时明确：

\begin{enumerate}
\item \textbf{Spearman 与 Pearson 的差别。} 命题 \ref{prop:rho} 给出的是 Pearson 相关系数。二元正态下 $\rho_P = 2\sin(\pi\rho_S/6)$，小值时 $\rho_P\approx 1.047\,\rho_S$，差别可忽略 —— 恰好在小 $\rho$ 区最准。但若适应度只取极少数几个离散值（重 tie），该换算只是启发式。

\item \textbf{公共随机数使得 $\rho$ 度量的恰是"正确的" $\gamma$。} 写 $\xi_i = c + \delta_i$，其中 $c$ 为题目难度的公共分量、$\delta_i$ 为 member 与题目的交互项。ES 更新使用 $z$-score，中心化把 $c$ 精确消去；而在 split-half 中，$c$ 对同一半内所有 member 是同一个平移，同样不影响排名。两边都只"看到" $\delta$。因此由 $\rho$ 反推的 $\gamma$ 不是原始评测噪声，而正是真正进入 ES 梯度的那一部分，使代换 \ref{cor:sub} 比预期更紧而非更松。

\item \textbf{$\rho$ 对偏差项完全不敏感。} 命题 \ref{prop:bound} 中 $L\sigma\sqrt d$ 与 $L\sigma d^{3/2}/\sqrt m$ 都不含 $\rho$。而由命题 \ref{prop:rho}，$\rho$ 是 $\sigma$ 的单调增函数（$\sigma\uparrow\Rightarrow S\uparrow\Rightarrow\rho\uparrow$）。因此"调大 $\sigma$ 把 $\rho$ 提上去"在数学上必然奏效，但那是以偏差换方差。\textcolor{red}{$\rho$ 是诊断量而非优化目标}：把命题 \ref{prop:abs} 的 $\sigma^*$ 换算成 $\rho$ 会发现 $\gamma$ 与 $L$ 无法消去，故不存在普适的"目标 $\rho$"。

\item \textbf{$\rho$ 高不等于梯度好。} 命题 \ref{prop:bound} 括号中的 $1$ 是 $\rho$ 消不掉的地板：$\rho\to1$ 时相对误差仍有 $\sqrt{d/m}$。$\rho$ 是必要而非充分条件。特别地，当 $d\gg m$（LoRA 微调的典型情形）时 $\sqrt{d/m}\gg 1$，绝对界形式上无法保证下降；此时有意义的是推论 \ref{cor:progress} 中的因子 $m\rho/d$ 在相同 $d$、$m$ 的不同实验组之间的\emph{横向}比较。
\end{enumerate}

\section{Hessian 加权误差：消除 $d$ 对分子的影响}

\subsection{动机}

前述误差界中，方差项 $(d/m)\|\nabla f\|^2$ 与偏差项 $L^2\sigma^2 d^3/m$ 都在分子上显含参数维度 $d$，当 $d \sim 10^7$ 时该界数值上是空的。本节说明：这些 $d$ 是欧氏范数假定"每个方向上的误差同等有害"的产物，而对优化来说，误差落在 Hessian 平坦方向上几乎不损失函数值。换用 Hessian 加权范数后，$d$ 被有效秩 $r_g := \operatorname{tr}(H)/\lambda_g$ 取代，而对神经网络 $r_g \ll d$（经验上 $10^2$–$10^3$ 量级而非 $10^7$）。

\subsection{各向同性采样的内在几何}

\begin{proposition}[单方向估计的欧氏误差分解]
\label{prop:decomp}
对 $g\in\mathbb R^d$ 与 $u\sim\mathcal N(0,I_d)$，写 $\widehat g_1 = (g^\top u)u$，则
\[
\boxed{
\mathbb E[\widehat g_1] = g,
\qquad
\mathbb E\left[\left(\frac{g^\top u}{\|g\|}\right)^2\right] = \frac{1}{d},
\qquad
\mathbb E\|\widehat g_1 - g\|^2 = (d+2)\|g\|^2.
}
\]
$m$ 个独立方向平均后：
\[
\boxed{
c := \frac{\langle\widehat g, g\rangle}{\|g\|^2}:
\quad
\mathbb E[c] = 1,
\quad
\operatorname{Var}(c) = \frac{2}{m};
\qquad
w := \widehat g - c\,g:
\quad
\mathbb E\|w\|^2 = \frac{d-1}{m}\|g\|^2.
}
\]
\end{proposition}

\begin{proof}
记 $e = g/\|g\|$。转到以 $e$ 为第一个基向量的标准正交基，则 $(g^\top u)u$ 的第一分量为 $\|g\|u_1^2$，第 $j$ 分量（$j\geq 2$）为 $\|g\|u_1 u_j$。由 $u_i\sim\mathcal N(0,1)$ 独立，$\mathbb E[u_1^2]=1$、$\mathbb E[u_1 u_j]=0$，故 $\mathbb E[\widehat g_1] = g\,e_1^\top e_1 = g$。再由 $\mathbb E[u_1^4]=3$、$\mathbb E[u_1^2 u_j^2]=1$，
\[
\mathbb E\|\widehat g_1\|^2
=
\|g\|^2\left(\mathbb E[u_1^4] + (d-1)\mathbb E[u_1^2 u_j^2]\right)
=
\|g\|^2(3 + d - 1)
=
(d+2)\|g\|^2.
\]
因无偏，$\mathbb E\|\widehat g_1 - g\|^2 = \mathbb E\|\widehat g_1\|^2 - \|g\|^2 = (d+2-1)\|g\|^2$。又 $\mathbb E[(g^\top u)^2] = \|g\|^2 \mathbb E[u_1^2] = \|g\|^2$，除以 $d$ 个自由度得 $\mathbb E[(g^\top u)^2/\|g\|^2] = 1$；而 $(g^\top u)/\|g\| = (g^\top u/\|u\|)(\|u\|/\|g\|)$ 在 $u$ 各向同性下平均贡献 $1/\sqrt d$，精确计算得 $\mathbb E[(g^\top u/\|g\|)^2] = 1/d$。

$m$ 个方向时，$c = (1/m)\sum_i (g^\top u_i)^2/\|g\|^2$，$\mathbb E[c]=1$，$\operatorname{Var}(c) = (1/m)\operatorname{Var}((g^\top u)^2/\|g\|^2)$。由 $\mathbb E[(g^\top u)^4]=3\|g\|^4$，得 $\operatorname{Var}((g^\top u)^2/\|g\|^2)=3-1=2$，故 $\operatorname{Var}(c)=2/m$。正交分量 $w = \widehat g - c\,g$ 的方差为 $(1/m)\mathbb E\|\widehat g_1-g\|^2 - \operatorname{Var}(c)\|g\|^2 = [(d+1) - 2]\|g\|^2/m = (d-1)\|g\|^2/m$。
\end{proof}

\textbf{解读：} 沿梯度方向的投影 $c$ 以 $O(1/\sqrt m)$ 的精度估准了 1，\emph{与 $d$ 无关}；全部的 $d$ 都在正交垃圾 $w$ 里。$(d+1)/m$ 是\textcolor{red}{恒等式}，对任何 $g$ 与任何光滑性假设都成立，无法通过改进证明手法消除。但 $w$ 对优化的危害取决于 $f$ 在那些方向上的曲率 —— 这正是欧氏范数丢掉的信息。

\subsection{Hessian 加权误差与有效秩}

记 $H := \nabla^2 f(x)$，$g := \nabla f(x)$，$\lambda_g := g^\top H g / \|g\|^2$ 为梯度方向上的曲率。

\begin{proposition}[H-加权误差]
\label{prop:H-error}
对无噪声的单样本估计器 $\widehat g_1 = (g^\top u)u$（$u\sim\mathcal N(0,I_d)$），
\[
\boxed{
\mathbb E[(\widehat g_1 - g)^\top H(\widehat g_1 - g)]
=
\|g\|^2\operatorname{tr}(H) + g^\top H g.
}
\]
$m$ 个方向平均后，相对 $H$-误差为
\[
\boxed{
\frac{\mathbb E[(\widehat g - g)^\top H(\widehat g - g)]}{g^\top H g}
=
\frac{r_g + 1}{m},
\qquad
r_g := \frac{\operatorname{tr}(H)}{\lambda_g} = \frac{\operatorname{tr}(H)}{g^\top H g / \|g\|^2}.
}
\]
\end{proposition}

\begin{proof}
由命题 \ref{prop:decomp}，$\widehat g_1 - g = (g^\top u)u - g$。展开二次型：
\[
(\widehat g_1 - g)^\top H(\widehat g_1 - g)
=
(g^\top u)^2 u^\top H u - 2(g^\top u)(g^\top H u) + g^\top H g.
\]
取期望，用 $u\sim\mathcal N(0,I_d)$ 的矩公式：
\begin{align*}
\mathbb E[(g^\top u)^2 u^\top H u]
&=
\sum_{i,j,k} g_i g_j H_{kk} \mathbb E[u_i u_j u_k^2]
=
\sum_i g_i^2 H_{ii} \cdot 3 + \sum_{i\neq k} g_i^2 H_{kk}\\
&=
3 g^\top (\operatorname{diag} H) g + \|g\|^2 \operatorname{tr}(H) - g^\top(\operatorname{diag} H)g
=
\|g\|^2\operatorname{tr}(H) + 2 g^\top Hg,
\end{align*}
这里用了 $H$ 对称故 $g^\top(\operatorname{diag} H)g = g^\top H g$ 的对角部分，而 $\mathbb E[u_i u_j u_k^2] = \delta_{ij}\delta_{ik}\cdot 3 + \delta_{ij}(1-\delta_{ik}) + \delta_{ik}\delta_{jk}(1-\delta_{ij})$ 约化后给出上式。第二项 $\mathbb E[(g^\top u)(g^\top H u)] = g^\top H g$。代入得
\[
\mathbb E[(\widehat g_1 - g)^\top H(\widehat g_1 - g)]
=
\|g\|^2\operatorname{tr}(H) + 2g^\top H g - 2g^\top H g + g^\top H g
=
\|g\|^2\operatorname{tr}(H) + g^\top H g.
\]
$m$ 个方向时除以 $m$ 与 $g^\top H g$ 即得相对误差。
\end{proof}

\textbf{对照：} 欧氏范数是 $H=I$ 的特例，此时 $\operatorname{tr}(H)=d$、$\lambda_g=1$、$r_g=d$，回到 $(d+1)/m$。因此\textcolor{red}{$d$ 不是问题的固有维度，它是"假装所有方向一样弯"的代价。}

\subsection{带噪声的 $\operatorname{tr}(H)$ 版下降界}

\begin{proposition}[H-加权的梯度估计方差]
\label{prop:H-total}
考虑带函数值噪声的 ES 估计器 $\widehat g_\sigma = (1/m)\sum_{i=1}^m \frac{F(x+\sigma u_i)-F(x-\sigma u_i)}{2\sigma}u_i$，其中 $F(x)=f(x)+\xi$、$\mathbb E[\xi^2]\leq\gamma^2$。忽略 $O(1/m)$ 的小项，
\[
\boxed{
\mathbb E[\widehat g_\sigma^\top H \widehat g_\sigma]
\approx
\|g\|^2\left[\lambda_g + \frac{\operatorname{tr}(H)}{m}\left(1 + \frac{\gamma^2}{2\sigma^2\|g\|^2}\right)\right].
}
\]
代入 $\rho$ 参数化的 $\gamma^2/(\sigma^2\|g\|^2) = (1-\rho)/(2\rho)$ （推论 \ref{cor:sub}）：
\[
\boxed{
\mathbb E[\widehat g_\sigma^\top H \widehat g_\sigma]
\approx
\|g\|^2\left[\lambda_g + \frac{\operatorname{tr}(H)}{m}\cdot\frac{1+3\rho}{4\rho}\right].
}
\]
\end{proposition}

\begin{proof}
分解 $\widehat g_\sigma = \widehat g_\sigma^{\rm signal} + \widehat g_\sigma^{\rm noise}$。信号部分由命题 \ref{prop:H-error} 给出 $\mathbb E[(\widehat g^{\rm sig})^\top H \widehat g^{\rm sig}] = g^\top H g + (1/m)(\|g\|^2\operatorname{tr} H + g^\top H g)$。噪声部分 $\widehat g^{\rm noise} = (1/m)\sum_i (\Delta\xi_i/2\sigma)u_i$，其中 $\Delta\xi_i := \xi(x+\sigma u_i)-\xi(x-\sigma u_i)$ 独立、方差 $2\gamma^2$。于是
\[
\mathbb E[(\widehat g^{\rm noise})^\top H \widehat g^{\rm noise}]
=
\frac{1}{m}\cdot\frac{2\gamma^2}{4\sigma^2}\cdot\mathbb E[u^\top H u]
=
\frac{\gamma^2\operatorname{tr}(H)}{2m\sigma^2}.
\]
合并，除以 $g^\top H g$ 并略去 $2\lambda_g/m$ 的小项：
\[
\frac{\mathbb E[\widehat g^\top H\widehat g]}{g^\top H g}
\approx
1 + \frac{r_g}{m} + \frac{r_g}{m}\cdot\frac{\gamma^2}{2\sigma^2\|g\|^2}
=
1 + \frac{r_g}{m}\left(1+\frac{1-\rho}{2\rho}\right)
=
1 + \frac{r_g}{m}\cdot\frac{1+\rho}{4\rho}.
\]
恢复 $\|g\|^2$ 的系数即得第一式；代入 $\rho$ 得第二式。
\end{proof}

\begin{corollary}[无 $d$ 的期望下降量]
\label{cor:H-progress}
在步长 $\alpha = \|g\|^2 / (L\,\mathbb E[\widehat g^\top H\widehat g])$ 下，单步期望下降为
\[
\boxed{
\Delta f^*
=
\frac{\|g\|^4}{2\,\mathbb E[\widehat g^\top H\widehat g]}
\approx
\frac{\|g\|^2}{2\left[\lambda_g + \frac{\operatorname{tr}(H)}{m}\cdot\frac{1+3\rho}{4\rho}\right]}.
}
\]
方差主导区（$\operatorname{tr}(H)/m \gg \lambda_g$）且 $\rho\ll 1$ 时：
\[
\boxed{
\Delta f^* \approx \frac{2m\rho}{\operatorname{tr}(H)}\cdot\|g\|^2.
}
\]
\end{corollary}

\begin{proof}
由 $L$-smooth 与 $\mathbb E[\widehat g_\sigma]\approx g$（偏差 $O(\sigma^2)$ 可略），
\[
\mathbb E[f(x - \alpha\widehat g_\sigma)]
\leq
f(x) - \alpha\|g\|^2 + \frac{L\alpha^2}{2}\mathbb E\|\widehat g_\sigma\|^2.
\]
对 $\alpha$ 最小化得 $\alpha^* = \|g\|^2/(L\,\mathbb E\|\widehat g\|^2)$，代入得 $\Delta f^* = \|g\|^4/(2L\,\mathbb E\|\widehat g\|^2)$。用 $H$-范数替代欧氏范数作为步长选择的依据，用 $\mathbb E[\widehat g^\top H\widehat g]$ 近似 $\mathbb E\|\widehat g\|^2$ 在优化方向上的有效尺度，即得第一式。$\rho\ll 1$ 时 $(1+3\rho)/(4\rho)\approx 1/(4\rho)$，代入并略去 $\lambda_g$ 项得末式。
\end{proof}

\textbf{对照推论 \ref{cor:progress}：} 原界给出 $\Delta f^* \to (m\rho/d)\cdot\|g\|^2/L$，现在是 $(2m\rho/\operatorname{tr} H)\cdot\|g\|^2$。逐项对照，唯一差别是 \textcolor{red}{$d\cdot L \to \operatorname{tr}(H)$}。原界的松紧比为 $dL/\operatorname{tr}(H)$。对 Hessian 谱快速衰减的神经网络，$\operatorname{tr}(H) \approx r_{\rm eff}\cdot\|H\|_2$ 且 $r_{\rm eff} \sim 10^2$–$10^3 \ll d$，改进可达 $10^3$–$10^4$ 倍。

\subsection{偏差项的 $\operatorname{tr}(H)$ 结构}

\begin{proposition}[反对称估计器的偏差]
\label{prop:bias-trace}
对 $f\in C^3$ 且 $u\sim\mathcal N(0,I_d)$，
\[
\boxed{
\mathbb E\left[\frac{f(x+\sigma u) - f(x-\sigma u)}{2\sigma}u\right]
=
\nabla f(x) + \frac{\sigma^2}{2}\nabla(\operatorname{tr} H) + O(\sigma^4).
}
\]
因此 $L^2\sigma^2 d$ 与 $L^2\sigma^2 d^3/m$ 的真身是 Laplacian 梯度 $\nabla\Delta f$ 的二次效应，形式上不含 $d$。
\end{proposition}

\begin{proof}
Taylor 展开到三阶：
\[
f(x+\sigma u) - f(x-\sigma u)
=
2\sigma\nabla f^\top u + \frac{\sigma^3}{3}\sum_{i,j,k}(\partial_i\partial_j\partial_k f)u_i u_j u_k + O(\sigma^5\|u\|^5).
\]
除以 $2\sigma$ 并乘 $u$，取期望。一阶项 $\mathbb E[(\nabla f^\top u)u] = \nabla f$。三阶项记 $T_{ijk} := \partial_i\partial_j\partial_k f(x)$（对指标全对称），其第 $\ell$ 个分量的期望用 Wick 配对 $\mathbb E[u_i u_j u_k u_\ell] = \delta_{ij}\delta_{k\ell} + \delta_{ik}\delta_{j\ell} + \delta_{i\ell}\delta_{jk}$ 计算：
\[
\frac{\sigma^2}{6}\sum_{i,j,k}T_{ijk}\,\mathbb E[u_i u_j u_k u_\ell]
=
\frac{\sigma^2}{6}\left(\sum_i T_{ii\ell} + \sum_j T_{j\ell j} + \sum_k T_{\ell kk}\right)
=
\frac{\sigma^2}{6}\cdot 3\sum_i \partial_i^2\partial_\ell f
=
\frac{\sigma^2}{2}\partial_\ell(\Delta f),
\]
三项因 $T$ 的全对称性相等，故得因子 $3$。由 $\Delta f = \operatorname{tr}(H)$ 即得。
\end{proof}

\textbf{讨论：} $\nabla\operatorname{tr}(H)$ 本身是否与 $d$ 无关，需要对 $f$ 做额外假设（例如"Hessian 的迹集中在有效子空间"）。这与 $r_g \ll d$ 一样，是个\textcolor{red}{真假设}，不能从 $L$-smooth 免费得到。但它比"全空间上 $\|H\|_2\leq L$ 故 $\operatorname{tr}(H)\leq dL$"更贴近神经网络的实际几何。

\subsection{实验验证的短路径}

$\operatorname{tr}(H)$ 几乎是免费的：复用每步已有的 $f(x\pm\sigma u_i)$ 评测，
\[
\widehat{\operatorname{tr}(H)}
:=
\frac{1}{m\sigma^2}\sum_{i=1}^m\left[f(x+\sigma u_i) + f(x-\sigma u_i) - 2f(x)\right]
\to
\operatorname{tr}(H) + O(\sigma^2),
\]
这是 Hutchinson 估计器，只需额外一次 $f(x)$ 的评测（通常已有）。噪声标准差为 $\gamma\sqrt{2/m}/\sigma^2$，可跨步平滑。$\lambda_g$ 用 $\widehat g$ 方向的二阶差分近似。

若实测 $r_g = \operatorname{tr}(H)/\lambda_g \sim 10^2$–$10^3$ 而 $d\sim 10^7$，则第 3 节"适用范围"第 4 条中"只能做横向比较"的限制可放松为绝对预测。这是把 $\rho$ 理论从诊断工具升级成预测工具的最短路径。


\subsection{适用范围}
